\chapter{Fréchet 微分}
\label{chap:frechet-differential-solutions}

\begin{solution}{operator-norm-basic-properties}
若 \(A\ne0\)，存在 \(x\) 使 \(Ax\ne0\)，归一化后可知
\(\norm A_{\mathrm{op}}>0\)；若 \(A=0\)，则定义中的上确界等于零。齐次性来自
\(\norm{\lambda Ax}=|\lambda|\norm{Ax}\)。对 \(\norm x\le1\)，
\[
 \norm{(A+B)x}\le\norm{Ax}+\norm{Bx}
 \le\norm A_{\mathrm{op}}+\norm B_{\mathrm{op}},
\]
取上确界得到三角不等式。若 \(A,B\) 可复合，则
\[
 \norm{ABx}\le\norm A_{\mathrm{op}}\norm{Bx}
 \le\norm A_{\mathrm{op}}\norm B_{\mathrm{op}}\norm x.
\]
对 \(\norm x\le1\) 取上确界即得乘法估计。
\end{solution}

\begin{solution}{matrix-operator-entry-bound}
对第 \(i\) 个坐标使用 Cauchy--Schwarz 不等式：
\[
 |(Ax)_i|^2
 =\left|\sum_{j=1}^na_{ij}x_j\right|^2
 \le\left(\sum_{j=1}^na_{ij}^2\right)\norm x_2^2.
\]
再对 \(i\) 求和，得到
\[
 \norm{Ax}_2^2\le
 \left(\sum_{i=1}^m\sum_{j=1}^na_{ij}^2\right)\norm x_2^2.
\]
开平方即为所求；因而算子范数不超过矩阵的 Frobenius 范数。
\end{solution}

\begin{solution}{affine-map-frechet-derivative}
若 \(F(x)=Ax+b\)，则
\[
 F(a+h)-F(a)-Ah=0.
\]
余项恒为零，当然是 \(o(\norm h)\)，故 \(DF(a)=A\)。这里不依赖
基点 \(a\)。
\end{solution}

\begin{solution}{frechet-partial-directional}
写 \(f(a+h)=f(a)+Ah+r(h)\)，其中
\(\norm{r(h)}/\norm h\to0\)。固定 \(v\ne0\)，有
\[
 \frac{f(a+tv)-f(a)}t
 =Av+\frac{r(tv)}t,
 \qquad
 \left\|\frac{r(tv)}t\right\|
 =\norm v\frac{\norm{r(tv)}}{\norm{tv}}\to0.
\]
故 \(D_vf(a)=Av\)。零方向直接成立。取 \(v=e_j\) 即得
\(\partial_jf(a)=Ae_j\)。
\end{solution}

\begin{solution}{euclidean-norm-derivative}
固定 \(x\ne0\)。由
\[
 \norm{x+h}-\norm x
 =\frac{2\langle x,h\rangle+\norm h^2}{\norm{x+h}+\norm x}
\]
以及当 \(h\to0\) 时分母趋于 \(2\norm x\)，可得
\[
 \norm{x+h}-\norm x
 =\frac{\langle x,h\rangle}{\norm x}+O(\norm h^2).
\]
例如把上式与 \(\langle x,h\rangle/\norm x\) 相减，利用
\(|\norm{x+h}-\norm x|\le\norm h\) 即可把余项估为
\(C_x\norm h^2\)。因此
\[
 D\norm{\cdot}(x)h=\frac{\langle x,h\rangle}{\norm x}.
\]
在原点，对任意 \(v\ne0\)，方向差商为
\[
 \frac{\norm{tv}}t=\frac{|t|}{t}\norm v.
\]
它在 \(t\to0^+\) 与 \(t\to0^-\) 时分别趋于 \(\norm v\) 与
\(-\norm v\)。本章方向导数使用双侧极限，所以非零方向的方向导数
均不存在，原点因而不可微。只有把方向极限改为单侧极限时，才会得到
\(\norm v\)；那不是本题采用的定义。
\end{solution}

\begin{solution}{partial-derivatives-discontinuous-example}
令 \(g(0,0)=0\)。当 \(t\ne0\) 时
\[
 \frac{g(t,0)-g(0,0)}t=0,
 \qquad
 \frac{g(0,t)-g(0,0)}t=0,
\]
所以两个偏导数均为零。另一方面，沿 \((t,t)\to(0,0)\)，
\(g(t,t)=1/2\)，不趋于 \(g(0,0)\)。故函数不连续，因而更不可能
Fréchet 可微。
\end{solution}

\begin{solution}{directional-map-not-linear}
对 \(v=(a,b)\ne0\)，
\[
 \frac{q(ta,tb)}t
 =\frac{a^3}{a^2+b^2}.
\]
所以方向导数等于右端。取 \(u=(1,0)\)、\(v=(0,1)\)，有
\[
 D_uq(0)=1,
 \qquad D_vq(0)=0,
 \qquad D_{u+v}q(0)=\frac12.
\]
最后一个数不等于前两个之和，故方向映射不可加。
\end{solution}

\begin{solution}{zero-directional-not-continuous}
设方向 \(v=(a,b)\)。若 \(b=0\)，则 \(h(ta,0)=0\)。若 \(b\ne0\)，
\[
 \frac{h(ta,tb)}t
 =\frac{t a^3b}{t^4a^6+b^2}\to0.
\]
故所有方向导数都是零。取曲线 \((x,y)=(t,t^3)\)，则 \(t\ne0\) 时
\[
 h(t,t^3)=\frac{t^6}{2t^6}=\frac12.
\]
它不趋于原点函数值零，所以函数不连续。若它 Fréchet 可微，则由全部方向
导数可知候选导数只能是零映射，但沿这条曲线的余项甚至不趋于零，矛盾。
\end{solution}

\begin{solution}{frechet-uniqueness-sequence-proof}
若余项满足可微定义，把极限限制到任意上述序列即可。反之，若定义中的
极限不是零，则存在 \(\varepsilon_0>0\)，对每个 \(k\ge1\) 可选
\(0<\norm{h_k}<1/k\)、\(a+h_k\in U\)，使余项比不小于
\(\varepsilon_0\)，这与序列假设矛盾。

唯一性只需在这个“对每个序列”条件中选择一个能检出两映射差异的序列。
若 \(A\ne B\)，取 \(v\) 使 \((A-B)v\ne0\)，并令 \(h_k=v/k\)。两种
展开相减后有
\[
 \frac{\norm{(A-B)h_k}}{\norm{h_k}}
 \le
 \frac{\norm{r_A(h_k)}}{\norm{h_k}}
 +\frac{\norm{r_B(h_k)}}{\norm{h_k}}\to0.
\]
左端恒等于 \(\norm{(A-B)v}/\norm v>0\)，矛盾。因此 \(A=B\)。
\end{solution}

\begin{solution}{continuous-partials-proof-details}
令
\(a^{(j)}=a+(h_1,\ldots,h_j,0,\ldots,0)\)。对固定 \(j\)，函数
\[
 t\longmapsto f(a^{(j-1)}+te_j)
\]
在连接 $0$ 与 $h_j$ 的区间上逐分量可微。对每个分量 $f_\ell$ 使用一元
中值定理；若 $h_j=0$，增量恒为零，直接取
$\xi_{j\ell}=a^{(j-1)}$。因此总能取该段上的点 $\xi_{j\ell}$，使
\[
 f_\ell(a^{(j)})-f_\ell(a^{(j-1)})
 =h_j\partial_jf_\ell(\xi_{j\ell}).
\]
将各段相加后，减去线性项
$\sum_jh_j\partial_jf(a)$。所有中值点与 $a$ 的距离不超过 $\norm h$；
有限多个标量偏导在 $a$ 连续，故它们与在 $a$ 处取值的差可以由同一个
趋零量控制。再用 $\sum_j|h_j|\le\sqrt n\norm h$，便得到
$o(\norm h)$ 余项。证明只需偏导在邻域存在并在 $a$ 连续，不需要它们在
整个邻域连续。
\end{solution}

\begin{solution}{scalar-product-rule-frechet}
写
\(f(a+h)=f(a)+Ah+r(h)\)、
\(g(a+h)=g(a)+Bh+s(h)\)。相乘并减去
\(f(a)g(a)\)，一次项为
\[
 f(a)Bh+g(a)Ah.
\]
其余各项包含 \(r(h)\)、\(s(h)\) 或两个 \(O(\norm h)\) 增量的乘积，
所以都是 \(o(\norm h)\)。故
\(D(fg)(a)=f(a)Dg(a)+g(a)Df(a)\)。
\end{solution}

\begin{solution}{reciprocal-rule-frechet}
一元函数 \(R(t)=1/t\) 在 \(t=g(a)\ne0\) 处导数为
\(-1/g(a)^2\)。链式法则给出
\[
 D(1/g)(a)=-\frac{Dg(a)}{g(a)^2}.
\]
若 \(f\) 也可微，再与乘积法则组合，得到
\[
 D(f/g)(a)=\frac{g(a)Df(a)-f(a)Dg(a)}{g(a)^2}.
\]
分母非零还由连续性保证在 \(a\) 的某邻域内商有定义。
\end{solution}

\begin{solution}{inner-product-derivative}
内积是连续双线性映射。令增量分别为
\(Du(a)h+o(\norm h)\) 与 \(Dv(a)h+o(\norm h)\)，双线性乘积法则给出
\[
 D\langle u,v\rangle(a)h
 =\langle Du(a)h,v(a)\rangle
  +\langle u(a),Dv(a)h\rangle.
\]
两增量的内积为 \(O(\norm h^2)\)，进入余项。
\end{solution}

\begin{solution}{quadratic-form-derivative}
展开
\[
 q(x+h)-q(x)
 =h^{\mathsf T}Ax+x^{\mathsf T}Ah+h^{\mathsf T}Ah.
\]
最后一项为 \(O(\norm h^2)\)，故
\[
 Dq(x)h=x^{\mathsf T}(A+A^{\mathsf T})h
 =\langle(A+A^{\mathsf T})x,h\rangle.
\]
把 \(A\) 分解为对称部分与斜对称部分；斜对称矩阵 \(S\) 满足
\(x^{\mathsf T}Sx=0\)，所以只有对称部分影响 \(q\) 和其导数。
\end{solution}

\begin{solution}{normalization-map-derivative}
把 \(N(x)=x\,\norm x^{-1}\) 看作向量与标量的乘积。由范数求导及商法则，
\[
 DN(x)h
 =\frac h{\norm x}
  -\frac{x\langle x,h\rangle}{\norm x^3}.
\]
与 \(x\) 作内积得到
\[
 \langle x,DN(x)h\rangle
 =\frac{\langle x,h\rangle}{\norm x}
  -\frac{\norm x^2\langle x,h\rangle}{\norm x^3}=0.
\]
这是因为归一化映射的像位于单位球面，其一阶增量沿球面的切方向。
\end{solution}

\begin{solution}{chain-rule-polar-map}
直接计算
\[
 J_F(r,\theta)=
 \begin{pmatrix}
  \cos\theta&-r\sin\theta\\
  \sin\theta&r\cos\theta
 \end{pmatrix}.
\]
若 \(f\) 在 \(F(r,\theta)\) 可微，链式法则给出
\[
 \partial_r(f\circ F)
 =f_x\cos\theta+f_y\sin\theta,
\]
\[
 \partial_\theta(f\circ F)
 =-r f_x\sin\theta+r f_y\cos\theta,
\]
其中 \(f_x,f_y\) 均在 \((r\cos\theta,r\sin\theta)\) 处取值。
\end{solution}

\begin{solution}{differentiable-derivative-discontinuous}
对 \(t\ne0\)，
\[
 \phi'(t)=2t\sin(1/t^2)-\frac2t\cos(1/t^2).
\]
在零点，
\[
 \frac{\phi(t)-\phi(0)}t=t\sin(1/t^2)\to0,
\]
所以 $\phi'(0)=0$。取 $t_k=(2\pi k)^{-1/2}$（$k\ge1$），则
\[
 \phi'(t_k)=-\frac2{t_k},
\]
其绝对值趋于无穷。令 $F(x_1,\ldots,x_n)=\phi(x_1)$。由
$|F(x)|\le x_1^2\le\norm x^2$，它在原点的 Fréchet 导数为零；但第一
偏导在原点附近无界。
\end{solution}

\begin{solution}{vector-mean-value-failure}
有 \(F(2\pi)=F(0)=(1,0)\)，而
\(F'(t)=(-\sin t,\cos t)\) 的范数恒为一。若存在精确向量中值点 \(c\)，
则 \(0=2\pi F'(c)\)，不可能。范数估计的证明先选取终点差方向上的单位
向量 \(u\)，再对标量函数 \(\langle u,F(t)\rangle\) 使用一元中值定理；
它只要求一个标量投影在某点达到所需变化率，不要求所有分量共用同一中值点。
\end{solution}

\begin{solution}{bounded-jacobian-lipschitz}
若 \(x,y\in K\)，凸性保证线段 \([x,y]\subseteq K\)。向量值平均估计
直接给出
\[
 \norm{f(y)-f(x)}\le M\norm{y-x}.
\]
非凸时，导数界只在 \(K\) 上成立，连接两点的线段可能离开 \(K\)。例如
取 \(U=\R\)、\(K=\{-1,1\}\) 及
\[
 f(x)=50+75x-25x^3.
\]
有 \(f'(-1)=f'(1)=0\)，所以 \(K\) 上导数上界 \(M=0\)，但
\(f(-1)=0\)、\(f(1)=100\)。这不否定平均估计，因为线段内部的导数没有
被 \(M\) 控制。
\end{solution}

\begin{solution}{zero-derivative-constant-convex}
凸情形下，对任意 \(x,y\in U\)，向量值平均估计以 \(M=0\) 给出
\(f(x)=f(y)\)。若只假设道路连通，取连接 \(x,y\) 的道路 \(\gamma\)。
每个 \(z\in U\) 都有半径 \(r_z>0\) 使
\(\overline B(z,r_z)\subset U\)，在该球内同一估计说明 \(f\)
常值。因此集合
\[
 A=\{t\in[0,1]:f(\gamma(t))=f(\gamma(0))\}
\]
在 \([0,1]\) 中既开又闭，并且非空；区间连通性给出 \(A=[0,1]\)。故
\(f(x)=f(y)\)。
\end{solution}

\begin{solution}{local-lipschitz-c1}
若 \(f\in C^1\)，在 \(a\) 周围取闭球
\(\overline B(a,r)\subset U\)。连续函数
\(x\mapsto\norm{Df(x)}_{\mathrm{op}}\) 在该紧球上有界。对更小开球中的
两点，连接线段仍在闭球内，平均估计给出同一个 Lipschitz 常数。

这只是局部结论。函数 \(f(x)=x^2\) 在 \(\R\) 上属于 \(C^1\)，但若有
全局 Lipschitz 常数 \(L\)，则
\(|(x+1)^2-x^2|=2x+1\le L\) 对所有正 \(x\) 成立，矛盾。
\end{solution}

\begin{solution}{directional-linear-not-frechet}
正文函数
\[
 h(x,y)=
 \begin{cases}x^3y/(x^6+y^2),&(x,y)\ne(0,0),\\0,&(x,y)=(0,0)
 \end{cases}
\]
已经满足要求。沿任意固定方向的差商趋于零，所以方向导数族是零线性映射。
若 Fréchet 导数存在，只能为零。沿 \((t,t^3)\to0\)，
\[
 \frac{|h(t,t^3)-0|}{\sqrt{t^2+t^6}}
 =\frac{1/2}{|t|\sqrt{1+t^4}}\longrightarrow+\infty,
\]
不满足余项为 \(o(\norm{(x,y)})\) 的条件。因此它不 Fréchet 可微。
\end{solution}
